\documentclass[11pt]{article}
\usepackage[a4paper,margin=27mm]{geometry}
\usepackage[T1]{fontenc}
\usepackage{lmodern}
\usepackage{amsmath,amssymb,amsthm}
\usepackage[colorlinks=true,linkcolor=blue,citecolor=blue,urlcolor=blue]{hyperref}
\hypersetup{pdftitle={Downstream transport control: a literature resolution of AIM 412},pdfauthor={Kaiwen Jin}}
\newtheorem{proposition}{Proposition}
\newtheorem{lemma}[proposition]{Lemma}
\newtheorem{corollary}[proposition]{Corollary}
\newcommand{\norm}[1]{\left\|#1\right\|}
\title{Downstream transport control:\\a literature resolution of AIM 412}
\author{Kaiwen Jin\\\small\href{mailto:yui.ui.siki@gmail.com}{\texttt{yui.ui.siki@gmail.com}}}
\date{12 October 2026}

\begin{document}
\maketitle
\begin{abstract}
AIM Problem 412 asks whether the worst minimal boundary-control norm for a
one-dimensional transport--diffusion equation tends to zero as the viscosity
vanishes, whenever the control time exceeds the transport time $L/M$.
The adjoint construction of Koike and Laheurte gives a negative answer:
for positive speed $M$, the cost instead diverges for every $0<T<2L/M$.
We explain its application to the problem's $L^2$ boundary controls and
transposition solutions, including the scaling to arbitrary $L,M>0$.
\end{abstract}

Fix $L,M>0$. The equation in Problem 412~\cite{aim} is
\begin{equation}\label{eq:primal}
\begin{cases}
y_t+My_x-\varepsilon y_{xx}=0 &(t,x)\in(0,T)\times(0,L),\\
y(t,0)=h(t),\quad y(t,L)=0 &t\in(0,T),\\
y(0,\cdot)=y_0\in L^2(0,L).
\end{cases}
\end{equation}
For $h\in L^2(0,T)$, the solution and its terminal nullity are interpreted
by transposition. Write
\begin{equation}\label{eq:cost}
C(T,\varepsilon)
=\sup_{\norm{y_0}_{L^2(0,L)}\leq1}
\inf\bigl\{\norm{h}_{L^2(0,T)}:\ y(T)=0\bigr\}.
\end{equation}
The infimum runs over admissible boundary controls.
The question is whether this cost tends to zero as $\varepsilon\downarrow0$
for every $T>L/M$.

Koike and Laheurte's preprint~\cite{kl}, submitted on 28 September 2026,
contains the following contrary conclusion; see its Theorem 1.3.
This source appeared after the problem's stated literature check of
22 September 2026. We give the direct implication of its adjoint theorem
for the precise formulation~\eqref{eq:primal}--\eqref{eq:cost}.

\begin{proposition}\label{prop:divergence}
For every fixed $L,M>0$ and every fixed $0<T<2L/M$,
\[
\lim_{\varepsilon\downarrow0}C(T,\varepsilon)=+\infty.
\]
\end{proposition}

The obstruction is a family of adjoint solutions with a nonvanishing
terminal interior norm and an observed boundary flux tending to zero.
The initial state in the worst-data cost may depend on $\varepsilon$;
this dependence is permitted by the supremum in~\eqref{eq:cost}.

\begin{lemma}[Koike--Laheurte]\label{lem:witness}
For each $0<\tau<2$, there are $c>0$ and real solutions $z_\delta$,
$\delta>0$, of
\[
(z_\delta)_s=(z_\delta)_\xi+\delta(z_\delta)_{\xi\xi},
\qquad z_\delta(s,0)=z_\delta(s,1)=0,
\]
on $(0,\tau)\times(0,1)$, with
$z_\delta(0,\cdot)\in H^2(0,1)\cap H_0^1(0,1)$, such that
\begin{equation}\label{eq:witness}
\norm{z_\delta(\tau,\cdot)}_{L^2(0,1)}^2\geq c,
\qquad
\int_0^\tau|\delta(z_\delta)_\xi(s,0)|^2\,ds\longrightarrow0
\quad(\delta\downarrow0).
\end{equation}
The lower bound holds for all sufficiently small $\delta>0$, with $c$
independent of $\delta$. The solutions constructed in the source proof
are smooth finite eigenfunction sums.
\end{lemma}

This is Theorem 4.1 of~\cite{kl}, stated on PDF page 16 and proved on
pages 16--18. Its construction is entirely analytic.

\begin{proof}[Proof of Proposition~\ref{prop:divergence}]
Set $\tau=MT/L$ and $\delta=\varepsilon/(ML)$, so $0<\tau<2$.
Using Lemma~\ref{lem:witness}, define
\[
Z_\varepsilon(t,x)=L^{-1/2}z_\delta(Mt/L,x/L).
\]
Differentiation gives
\[
(Z_\varepsilon)_t=ML^{-3/2}(z_\delta)_s,
\quad (Z_\varepsilon)_x=L^{-3/2}(z_\delta)_\xi,
\quad (Z_\varepsilon)_{xx}=L^{-5/2}(z_\delta)_{\xi\xi}.
\]
Thus $Z_\varepsilon$ solves the forward adjoint equation
$(Z_\varepsilon)_t=M(Z_\varepsilon)_x+
\varepsilon(Z_\varepsilon)_{xx}$ with homogeneous Dirichlet endpoints.
Changing variables in the terminal and observation norms yields
\begin{align}
a_\varepsilon^2
&:=\norm{Z_\varepsilon(T,\cdot)}_{L^2(0,L)}^2
=\norm{z_\delta(\tau,\cdot)}_{L^2(0,1)}^2\geq c,
\label{eq:terminal}\\
b_\varepsilon^2
&:=\int_0^T|\varepsilon(Z_\varepsilon)_x(t,0)|^2\,dt
=M\int_0^\tau|\delta(z_\delta)_\xi(s,0)|^2\,ds
\longrightarrow0.\label{eq:flux}
\end{align}
Here $L,M,T$ are fixed, and $\delta\downarrow0$ precisely when
$\varepsilon\downarrow0$.

Put $p(t,x)=Z_\varepsilon(T-t,x)$. It solves
$p_t+Mp_x+\varepsilon p_{xx}=0$ with zero endpoint traces.
For smooth primal solutions, integration by parts gives
\[
\frac{d}{dt}\int_0^L yp\,dx
=\bigl[-Myp+\varepsilon(y_xp-yp_x)\bigr]_0^L
=\varepsilon h(t)p_x(t,0).
\]
The corresponding terminal identity is
\begin{equation}\label{eq:duality}
\langle y(T),Z_\varepsilon(0)\rangle
-\langle y_0,Z_\varepsilon(T)\rangle
=\int_0^T h(t)\varepsilon(Z_\varepsilon)_x(T-t,0)\,dt.
\end{equation}
For $L^2$ boundary inputs, this is the defining terminal duality identity
of the transposition solution. The first pairing is a terminal duality
pairing, which vanishes when $y(T)=0$; no stronger terminal trace is needed.
The smooth adjoint test makes the boundary integral well-defined.

Choose $y_{0,\varepsilon}=Z_\varepsilon(T)/a_\varepsilon$.
By~\eqref{eq:terminal}, this is an admissible unit $L^2$ initial state.
Every null control for it satisfies, by~\eqref{eq:duality} and
Cauchy--Schwarz,
\begin{equation}\label{eq:lower}
a_\varepsilon
\leq\norm{h}_{L^2(0,T)}b_\varepsilon.
\end{equation}
Time reversal leaves the observation norm $b_\varepsilon$ unchanged.
If $b_\varepsilon>0$, taking the infimum for this initial state and then
the supremum in~\eqref{eq:cost} gives
$C(T,\varepsilon)\geq a_\varepsilon/b_\varepsilon$.
If $b_\varepsilon=0$,~\eqref{eq:lower} rules out a null control for this
state; with $\inf\varnothing=+\infty$, the same conclusion of divergence
holds. Since $a_\varepsilon\geq\sqrt c$ for all sufficiently small
$\varepsilon$ and $b_\varepsilon\to0$, this proves the proposition.
No minimizer or interchange of the supremum with the limit is used.
\end{proof}

The source leaves the critical time $T=2L/M$ open~\cite[PDF p.~5]{kl}.
That endpoint is unnecessary for the original yes/no question:

\begin{corollary}[Negative answer to AIM 412]
For every fixed $L,M>0$, the time $T_0=3L/(2M)$ satisfies $T_0>L/M$ and
\[
\lim_{\varepsilon\downarrow0}C(T_0,\varepsilon)=+\infty.
\]
Consequently the assertion that this cost tends to zero for every
$T>L/M$ is false.
\end{corollary}
Here $L$ is the interval length, $M$ is the positive transport speed,
$\varepsilon$ is the viscosity, and $C$ is the worst minimal $L^2$
boundary-control norm in~\eqref{eq:cost}.
The corollary is the specialization $T=3L/(2M)$ of
Proposition~\ref{prop:divergence}.

\begin{thebibliography}{2}
\bibitem{kl}
K. Koike and V. Laheurte,
\emph{Uniform null-controllability times for the Coron--Guerrero problems
are $2$ and $2+2\sqrt2$},
arXiv:2609.35355v1, submitted 28 September 2026.
\url{https://arxiv.org/abs/2609.35355v1}.

\bibitem{aim}
M. Colbrook's AIM problem collection,
\emph{412. Sharp vanishing-viscosity control time for downstream transport},
revision \texttt{c9929805cece044705abaaebe1cebddac1ba11a6}.
\href{https://github.com/MColbrook/AIM/blob/c9929805cece044705abaaebe1cebddac1ba11a6/problems/412-positive-transport-sharp-uniform-control-time.md}{Immutable problem statement}.
\end{thebibliography}

\paragraph{AI assistance.}
OpenAI Codex was used to retrieve the cited sources, check the connection
between the adjoint theorem and the stated problem, and draft this note.
The mathematical threshold result is attributed to Koike and Laheurte.
No independent human review is asserted; independent mathematical review
of this literature-resolution report is requested.
\end{document}
